\documentclass{article}
\usepackage[UTF8]{ctex}
\usepackage{amsmath,amssymb,amsfonts}
\usepackage{enumitem}
\usepackage{geometry}
\geometry{a4paper, margin=1in}
% === 错题本排版优化（不参与题目标识） ===
\setlength{\parindent}{0pt}
\setlength{\parskip}{0.35em plus 0.10em minus 0.08em}
\setlength{\emergencystretch}{3em}
\setlist{itemsep=0.25em, topsep=0.40em, parsep=0pt}
\allowdisplaybreaks[2]
% === 错题本排版优化结束 ===
\begin{document}
\section*{南开大学2024年高等代数真题}
\section*{48．南开大学2024年高等代数试卷}

\subsection{题目1}
1．（30 分）若

\[
f(x)=x^{5}-x^{3}-x-1, g(x)=x^{4}+2 x^{2}+5 x+2 .
\]

计算 \(f(x)\) 与 \(g(x)\) 的首项系数为 1 的最大公因式 \((f(x), g(x))\) ．

\subsection{题目2}
2．（20 分）计算行列式

\[
\left|\begin{array}{cccc}
1+x & x & 0 & 0 \\
x & x+x^{2} & x^{2} & 0 \\
0 & x^{2} & x^{2}+x^{3} & x^{3} \\
0 & 0 & x^{3} & x^{3}+x^{4}
\end{array}\right|
\]

\subsection{题目3}
3．（20分）若四阶方阵 \(A\) 的伴随矩阵为

\[
\left(\begin{array}{cccc}
-1 & -1 & 1 & 0 \\
0 & -1 & 0 & 1 \\
0 & 0 & -1 & -1 \\
0 & 0 & 0 & -1
\end{array}\right)
\]

求 \(A^{3}\) ．

\subsection{题目4}
4．（20 分）设 \(V=\mathbb{R}[x]_{4}\) 为次数小于 4 的实系数多项式构成的实线性空间，定义线性变换 \(\mathcal{A}: V \rightarrow V\)如下：

\[
\mathcal{A}(f(x))=\left(1-x^{2}\right) f^{\prime \prime}(x)-2 x f^{\prime}(x), \forall f(x) \in V .
\]

求 \(V\) 的一组基，设 \(\mathcal{A}\) 在这组基下的矩阵为对角矩阵。

\subsection{题目5}
5．（20 分）求矩阵

\[
A=\left(\begin{array}{cccc}
1 & 1 & 0 & 2 \\
0 & 2 & -1 & 0 \\
0 & 1 & 0 & 0 \\
0 & -1 & 1 & 1
\end{array}\right)
\]

的若尔当标准形。

\subsection{题目6}
6．（20 分）设 \(\widetilde{A}, \widetilde{B} \in \mathbb{R}^{m \times(n+1)}\) 分别是两个线性方程组的增广矩阵，如果两个线性方程组都有解且解集相同，证明：\(\widetilde{A}, \widetilde{B}\) 的行向量组等价。

\subsection{题目7}
7．（10 分）已知 \(A, B \in \mathbb{R}^{n \times n}\) ，且 \(A\) 为正定矩阵，\(B\) 为半正定矩阵，证明：\(|A+B| \geq|A|\) 。

\subsection{题目8}
8．（10 分）设 \(V=\mathbb{C}^{n \times n}\) 为 \(n\) 阶复方阵构成的复线性空间，\(A \in V\) ，线性变换 \(\mathcal{D}: V \rightarrow V\) 定义为

\[
\mathcal{D}(X)=A X-X A^{\prime}, \forall X \in V .
\]

如果 \(A\) 有 \(n\) 个互不相同的特征值，证明：\(V=\mathcal{D}^{-1}(0) \oplus \mathcal{D}(V)\) ．
\end{document}
